\section{Method}
\label{sec:method}

\subsection{System, mission and invariants}

The converter is a non-ideal asynchronous Buck converter that is
\emph{physically of integer order}, supplied by a fixed source
$V_\mathrm{in}=48$~V, with $L=100~\mu$H, $C=100~\mu$F, $R=10~\Omega$,
$r_L=0.1~\Omega$, $r_C=0.05~\Omega$, $R_\mathrm{on}=0.05~\Omega$ and a
freewheeling diode with affine law $v_D=V_f+r_d\,i_D$, $V_f=0.42$~V,
$r_d=0.020~\Omega$. Fractional order belongs exclusively to the controller.

The imposed mission is $24\to1\to46$~V, with transitions at $0.108$ and
$0.216$~s and end at $0.324$~s, a $+0.5$~A load step between $0.054$ and
$0.090$~s, and initial conditions $(i_L,v_C)=(2.4~\mathrm{A},24~\mathrm{V})$
in all circumstances. The controller is sampled at $T_s=10~\mu$s with a
one-sample delay; PWM runs at $50$~kHz and the duty ratio is bounded to
$[0.02;0.98]$. The nominal feedforward
\begin{equation}
d_\mathrm{ff}=\frac{(R+r_L+r_d)\,v_\mathrm{ref}}{R\,V_\mathrm{in}-(R_\mathrm{on}-r_d)\,v_\mathrm{ref}}
\label{eq:dff}
\end{equation}
is kept without a $V_f$ term, with nominal variables frozen during
disturbances.

\paragraph{Cost and constraints.}
\begin{equation}
J=\frac{1}{T}\!\int_0^T\!\frac{|r-v_o|}{\max(r,1)}\,dt
+0.01\,\frac{\int_0^T i_L^2\,dt}{T\cdot 20^2}
+0.001\,\frac{\overline{|\Delta d|}}{0.96},\qquad T=0.324~\mathrm{s},
\end{equation}
subject to six constraints $g_i\le0$:
$g_1=I_\mathrm{max}/20-1$,
$g_2=(V_\mathrm{max}-46)/1.5-1$,
$g_3=|e_{24}|/0.05-1$,
$g_4=|e_{1}|/0.03-1$,
$g_5=|e_{46}|/0.10-1$,
$g_6=(0.8-V_\mathrm{min})/0.8$.
Extrema are taken on raw solver outputs and the errors $e_\bullet$ are means
over the steady-state windows $[0.0936;0.1044]$, $[0.198;0.2124]$ and
$[0.306;0.3204]$~s. The \emph{margin} of a candidate is $m=-\max_i g_i$, and
over a set of cases $m=\min_k(-\max_i g_{i,k})$: a minimum of maxima, hence a
non-differentiable function.

\subsection{Switched evaluator}

\paragraph{Why a switched model.}
The workspace of the reference Simulink model states that the 1~V plateau
lies within $3~\%$ of the boundary between continuous and discontinuous
conduction, and that the $24\to1$~V step forces about $3.2$~ms of
discontinuous conduction. An averaged continuous-conduction evaluator would
therefore be wrong exactly where $g_4$ applies. We simulate the switched
converter with three conduction states: switch on, diode on, and diode
blocked ($i_L=0$), the latter producing discontinuous conduction without a
separate model.

\paragraph{PWM edges.}
Each integration step (fourth-order Runge--Kutta) is split exactly at carrier
edges. Without this, a $1~\mu$s step over a $20~\mu$s period would limit
duty-ratio resolution to $5~\%$, making the $[0.02;0.98]$ bounds fictitious.
The relative difference in $J$ between the search step ($1~\mu$s) and the
confirmation step ($0.2~\mu$s) is \ChiffreEcartPas{} for a reference control
law.

\paragraph{Fractional realisation.}
The operator $s^\gamma$, $\gamma\in[0;1]$, is approximated by a
half-order-3 Oustaloup filter (7 cells) over $[1;10^5]$~rad/s, each cell
$(s+z)/(s+p)$ being discretised exactly with a zero-order hold. The integral
is realised as $s^{-\lambda}=s^{-1}\cdot s^{1-\lambda}$: an exact integrator
ensures zero steady-state error, and the Oustaloup cell carries only the
fractional part. The derivative uses the same realisation, including for
$\mu=1$, where the Oustaloup formula yields $\omega_h(s+\omega_b)/(s+\omega_h)$,
a band-limited derivative.

\subsection{FO-LQI synthesis and projection onto FO-PIDF}

The small-signal model is normalised by $I=5$~A, $V=24$~V, $d=0.48$ and
$T_0=1$~ms, then augmented with a memory state. The LQR law ($R=1$) gives
$K=[K_i^x,K_v^x,K_z]$. For the nominal linear plant, without disturbance and
without saturation, the state feedback reduces algebraically to
\begin{equation}
C_\mathrm{FOLQI}(s)=K_v^x+\frac{a}{R}
+C\,(a-r_CK_v^x)\,\frac{N s}{s+N}
+\frac{k_i}{s^\alpha},\qquad
a=K_i^x\frac{V}{I},\quad N=\frac{1}{r_CC},\quad k_i=-\frac{K_z}{T_0^{\alpha}},
\label{eq:projection}
\end{equation}
which gives the anchor $k_p=K_v^x+a/R$, $k_d=C(a-r_CK_v^x)$,
$\omega_f=\mathrm{clip}(N;20;2\cdot10^4)$, $b=1$, $c=0$. The factor
$T_0^{\alpha}$ is the path through which the parent order $\alpha$ acts on the
synthesis: the memory state is an integral of order $\alpha$ in normalised
time. The integral order of the projection is then replaced by $\lambda$,
and the derivative order is $\mu$.

In the integer-order case ($\alpha=\lambda=\mu=1$, exact pole),
relation~\eqref{eq:projection} reproduces proposition~1 of the earlier study
with a relative error of \ChiffreErreurPropUn{} over $[1;10^6]$~rad/s.
\textbf{This is the only exact case.} Projecting the pole
$N=2\cdot10^5$~rad/s into the band $[20;2\cdot10^4]$~rad/s alone introduces
a relative error of \ChiffreErreurProjPole{} at the top of the band. The
FO-LQI $\to$ FO-PIDF identity is therefore neither general nor exact; we
treat it as an anchor, never as an equivalence nor as a proof of fractional
optimality.

\subsection{Bound-saturation audit}

An earlier campaign spent more than 1800 evaluations without reaching robust
feasibility, and its diagnosis blamed clipping of $\log_{10}k_p$ at the bound
$-4$, reached by $79.0~\%$ of candidates. We therefore treat each bound as an
object of measurement: for each component of $\theta$, we report the
fraction of projections that reach each bound, over \ChiffreAuditN{} Sobol
samples of the decision box. We distinguish two kinds of saturation.
\emph{Structural} saturation affects a component that is constant in the
synthesis: the bound then removes no degree of freedom, it merely makes an
identity visible. \emph{Pathological} saturation affects a component that
varies and runs into the bound: the bound then truly truncates the search.

\subsection{Staged search and six metaheuristics}

\paragraph{Decision vector.}
The v5 vector acts on the \emph{synthesis}, not on the final controller:
$x=[\log q_i,\log q_v,\log q_\mathrm{mem},\alpha,\mu,\log_{10}\tau_\mathrm{aw},\lambda]$,
with $\lambda$ in seventh position and $\mu$ in fifth.

\paragraph{Methods.}
PSO, GA, ABC and three project variants --- pAEABC, pIGWO and pIGWO-DLH ---
reuse the operators of the earlier study. The prefix ``p'' signals that
operator-level identity with the original publications is not established.
All receive the same box, the same initial population for a given seed
(population 10), the same Deb comparator and the same budget counted in
\emph{evaluator calls}; pIGWO-DLH spends two calls per comparison and thus
performs fewer iterations at equal budget, by design.

\paragraph{Comparator.}
Feasible precedes infeasible; among infeasible points, the lower violation
wins; among feasible points, we rank by \emph{decreasing margin} rather than
increasing $J$. Ranking by $J$ produces candidates stuck to the boundary
(median nominal margin $0.0066$ in the earlier study, robust margin/violation
correlation $-0.821$), which are the worst starting points for robustness.
The final criterion is lexicographic: feasibility on all design cases, then
maximal minimum margin, then worst $J$ at equal margin ($<10^{-4}$).

\paragraph{Stages.}
Stage~1 seeks nominal feasibility on a single scenario, with
\ChiffreBudgetUn{} calls; a seed producing fewer than \ChiffreMinFaisables{}
distinct feasible points is declared \emph{sterile} and reported as such.
Stage~2 seeks robustness on the design set with \ChiffreBudgetDeux{} calls,
seeded only from the stage-1 feasible points of the \emph{same} seed.
Stage~3 replays every retained candidate at a $0.2~\mu$s step: no candidate is
declared feasible without this replay.

\subsection{Validation protocol}

\paragraph{Uncertainty set (declared choice).}
The brief mentions seventeen robust cases and sixteen corners without
defining them. We retain $L$, $C$ and $R$ at $\pm20~\%$ and $r_C$ at
$\pm50~\%$, i.e.\ $2^4=16$ corners plus the nominal case. $r_C$ is chosen
because it alone sets the projected pole $\omega_f=1/(r_CC)$, hence the
projection error. $V_\mathrm{in}$ stays fixed and the mission is never
modified.

\paragraph{Three disjoint sets.}
The sixteen corners are split between design and validation by random draw
(published seed \ChiffreGraineSplit), never after observing which cases fail.
The test set consists of uniform draws in the hypercube, with seed
\ChiffreGraineTest{} sealed in the protocol before any evaluation, and is
played only once; a sentinel file prevents replaying it. A set consulted for a
decision becomes a design set and loses all probative value.

\paragraph{Test size.}
Success on $n$ draws without failure certifies nothing beyond the 95~\% Wilson
interval. With $50/50$, the lower bound is $92.9~\%$. The brief stated that
``about 300'' draws suffice for a $99~\%$ lower bound; the exact computation
gives $98.74~\%$ for $300/300$ and requires $381$ draws. We retain
\ChiffreNTest.

\subsection{Preregistered statistics}

The statistical unit is the seed: thirty common seeds, paired methods. We
apply the Friedman omnibus test (with the Iman--Davenport correction), then
pairwise Wilcoxon signed-rank tests corrected with Holm's method, and report
the paired rank-biserial correlation as effect size. Success proportions are
given with their 95~\% Wilson interval. All seeds are reported, sterile ones
included; no threshold is adjusted after observation.
